\chapter{Convergence Insufficiency}
\index{Convergence Insufficiency}
\label{sec:Convergence Insufficiency}

\section{Overview}
\begin{itemize}[noitemsep]
  \item Convergence insufficiency is an eye coordination disorder in which the eyes fail to turn inward sufficiently when looking at a near object
  \item It affects roughly 5\% of children and adults and is a common cause of eye strain and reading difficulty
  \item It is not a vision-threatening condition but can significantly affect schoolwork and daily near tasks
\end{itemize}
\section{Causes}
This condition happens because of impaired coordination between the eye muscles that control inward turning (convergence). The underlying cause is often unclear; it may be related to the brain's ability to control binocular alignment. It is more common after head injury or concussion, and can coexist with other binocular vision problems.
\section{Risk Factors}
\begin{itemize}[noitemsep]
  \item Family history of eye coordination problems
  \item Head trauma or concussion
  \item Prolonged near work (reading, screens)
  \item Neurodevelopmental conditions such as attention-deficit/hyperactivity disorder
  \item Fatigue and stress
\end{itemize}
\section{Signs and Symptoms}
\subsection{Common symptoms}
\begin{itemize}[noitemsep]
  \item Eye strain, fatigue, and headache with near work
  \item Blurred or double vision when reading
  \item Words appear to move or "swim" on the page
  \item Squinting, rubbing eyes, or closing one eye while reading
  \item Poor concentration and reading comprehension, avoiding near tasks
\end{itemize}
\subsection{Emergency warning signs}
\begin{itemize}[noitemsep]
  \item Sudden onset of double vision or headache after head injury requires urgent medical evaluation
\end{itemize}
\section{Progression and Stages}
Symptoms typically appear in childhood or adolescence and are often first noticed when school reading demands increase. The condition tends to persist if untreated. It does not worsen vision or damage the eyes, but symptoms can fluctuate with fatigue, stress, and the amount of near work.
\section{Complications}
\begin{itemize}[noitemsep]
  \item Academic underperformance and avoidance of reading
  \item Chronic headaches and eye strain
  \item Reduced attention and self-esteem in children
  \item Difficulty with prolonged near work in adults (screen-based occupations)
\end{itemize}
\section{Diagnosis}
\begin{itemize}[noitemsep]
  \item Comprehensive eye examination with assessment of binocular function
  \item Near point of convergence testing
  \item Cover test to detect eye misalignment
  \item Measurement of convergence amplitudes
  \item Evaluation of accommodative function
\end{itemize}
\section{Prevention}
\begin{itemize}[noitemsep]
  \item Convergence insufficiency is largely not preventable
  \item Regular eye examinations for children, especially with reading complaints
  \item Prompt evaluation after concussion or head injury
  \item Ergonomic practices for near work
\end{itemize}
\section{Treatment Options}
\begin{itemize}[noitemsep]
  \item Vision therapy (office-based orthoptic exercises) is the most effective treatment, improving convergence in most cases
  \item Home-based pencil push-up exercises as a first-line option
  \item Prismatic glasses (base-in) for symptom relief in some patients
  \item Regular monitoring and retraining of eye coordination
  \item Treatment of coexisting accommodative disorders
\end{itemize}
\section{Living with It}
\begin{itemize}[noitemsep]
  \item Practice prescribed eye exercises regularly and consistently
  \item Take frequent breaks from near work and screens
  \item Ensure good lighting and posture when reading
  \item Teachers and parents should be informed about accommodations
\end{itemize}
\section{Outlook}
The outlook is very good. With vision therapy, most people achieve normal convergence and resolution or significant improvement of symptoms. Untreated, the condition persists but does not damage the eyes or vision.
